The Segre variety is covered by the charts $z_{ij}\ne0$, identified with $U_i\times V_j\cong\mathbb A^n\times\mathbb A^m$ through the ratios $x_a/x_i=z_{aj}/z_{ij}$ and $y_b/y_j=z_{ib}/z_{ij}$. The two projections are therefore morphisms.

(b) If $X,Y$ are projective, $X\times Y$ is the intersection of the inverse images of these closed sets under the projections, so it is closed in the Segre variety. Its nonempty product affine charts are irreducible by \exref{I.3.15}, and any two intersect in a nonempty open subset, so it is irreducible and projective.

(a) In general, $X\times Y$ is open in $\overline X\times\overline Y$, namely the intersection of the inverse images of $X$ and $Y$. Hence it is quasi-projective and irreducible.

(c) Given morphisms $f:Z\to X$ and $g:Z\to Y$, the unique set map $z\mapsto(f(z),g(z))$ is a morphism on the inverse images of the product affine charts by \exref{I.3.15}(c). Thus it is a morphism globally, proving the universal property.
